\documentstyle[%


righttag%


,11pt%


,amssymb%


,russian%               remove in Eng.


%,izvru%                 Set izven% in Eng.


,izvru-ks%             for brief communications


]{amsart}





% ! \text{...}





\newcommand{\dom}{\operatorname{dom}}


\newcommand{\im}{\operatorname{im}}


\newcommand{\re}{\operatorname{Re}}


\newcommand{\tts}[1]{}





\theoremstyle{plain}


\theorembodyfont{\it}


\newtheorem{lemnonum}{\hskip\parindent Лемма}


\renewcommand{\thelemnonum}{}





%\hoffset=-15mm%                  For Eng.


\setcounter{page}{81}


\god{2005}


\nomer{4~(515)} %                        Remove in Eng.


%\nomer{Vol.\,49, No.\,4}%               Set in Eng.


%\str{\pageref{begin}--\pageref{end}}%  Set in Eng.


\udk{517.9}





\author{\it В.Е.\,Федоров, М.А.\,Сагадеева}


% NB:   ^^ remove in Eng.


\title{Об ограниченных на прямой решениях


линейных уравнений соболевского типа с относительно


секториальными операторами}


\thanks{Работа выполнена при частичной финансовой поддержке гранта


для молодых ученых Правительства


Челябинской области и гранта Министерства образования Российской Федерации


(\No\,PD02-1.1-82).}





\date{}


%\pagestyle{myheadings}%         Set in Eng.


\pagestyle{plain} %              Remove in Eng.


\begin{document}


%\thispagestyle{plain}%          Set in Eng.


\maketitle


\label{begin}








\section*{Введение}





Пусть ${\cal U}$ и ${\cal F}$ ---


банаховы пространства, оператор $L\in \cal L(\cal U;\cal F)$ линейный


непрерывный, $\ker L\ne\{0\}$, оператор $M\in \cal C l(\cal


U;\cal F)$ линейный замкнутый с плотной областью определения,


функция $f:{\Bbb R}\to{\cal F}$. Многие начально-краевые задачи


для уравнений и систем уравнений, моделирующих различные реальные


процессы ([1]; [2], с.\,13--15), сводятся к задаче Коши


\begin{equation}


u(0)=u_0


\end{equation}


для операторно-дифференциального уравнения соболевского типа [1]--[4]


\begin{equation}


L\Dot{u}(t)=Mu(t)+f(t).


\end{equation}


В данной работе изучаются ограниченные на всей прямой 


${\Bbb R}$ классические решения уравнения (2)


или соответствующей задачи Коши.





В случае существования оператора $L^{-1}\in\cal L(\cal F;\cal U)$


уравнение (2) можно привести к уравнению $\Dot


u(t)=Su(t)+w(t)$, где $S=L^{-1}M\in{\cal C}l(\cal U)$,


$\dom S=\dom M$, $w(t)=L^{-1}f(t):{\Bbb R}\to {\cal U}$. Условия


существования ограниченных решений уравнений такого вида с


оператором $S\in\cal L(\cal U)$ получены в ([5], с.\,118--131).





В [6]--[8] изучены аналогичные вопросы для


нестационарного уравнения


$$


\Dot u(t)=S(t)u(t)+w(t),\quad t\in J\subset \Bbb R\,.


$$


При этом в ([6], с.\,145--343)


исследован случай оператор-функции $S(t):J\to\cal L(\cal U)$. Более


общий случай изучен в


([7], c.\,165--181; [8], c.\,245--252), где значения оператор-функции $S(t)$


являются неограниченными операторами при $t\in J$. Кроме того, в


([6], с.\,145--343; [7], c.\,165--181) получены условия


существования ограниченных решений и из более общих классов.





Вопросы существования ограниченных решений уравнения (2) изучены


в [9], где было рассмотрено это уравнение с


$(L,\sigma)$-ограниченным и с сильно $(L,p)$-сек\-то\-ри\-аль\-ным


оператором $M$. В последнем случае речь шла только о решениях,


заданных на положительной полуоси.





Для исследования вопросов существования ограниченных на всей оси


решений уравнения (2) с сильно


$(L,p)$-секториальным неограниченным оператором $M$


применим методы, используемые в ([5], с.\,118--131 и [9]).


В работе получены достаточные условия существования единственного 


ограниченного на всей прямой решения уравнения (2) и задачи Коши для него. 


Кроме того, доказана теорема о необходимом условии существования единственного


ограниченного на ${\Bbb R}$ решения исследуемого уравнения.








\section{Относительно $p$-секториальные операторы}





Утверждения данного параграфа доказаны в [1], [10].





Введем обозначения


\begin{gather*}


\rho^L(M)=\{\mu\in{\Bbb C}:(\mu L-M)^{-1}


\in\cal L(\cal F;\cal U)\},\quad


\sigma^L(M)={\Bbb C}\setminus\rho^L(M), \\


%


R_{\mu}^L(M)=(\mu L-M)^{-1}L,\quad L_{\mu }^L(M)=L(\mu L-M)^{-1}, \\


%


R_{(\lambda ,p)}^L(M)=\prod


\limits_{k=0}^pR_{\lambda _k}^L(M),\quad


L_{(\lambda ,p)}^L(M)=\prod\limits_{k=0}^pL_{\lambda _k}^L(M).


\end{gather*}








\begin{defn}


Оператор $M$ называется {\it сильно


$(L,p)$-секториальным}, если существуют константы $K>0$, $a\in {\Bbb R}$,


$\theta\in ( \pi /2, \pi )$ такие, что сектор 


$S_{a,\theta }^L(M)=\{\mu\in {\Bbb C}:|\arg (\mu -a)|<\theta$,


$\mu\ne a\}\subset\rho^L(M)$, 


причем для всех $ \mu _k\in S_{a,\theta}^L(M)$, $k=\overline{0,p}$, 


$$


\max\{\|R_{(\mu ,p)}^L(M)\|_{\cal L(\cal U)},\| L_{(\mu


,p)}^L(M)\|_{\cal L(\cal F)}\}\le K\Big/\prod\limits


_{k=0}^p|\mu _k-a|,


$$


и существует плотный в ${\cal F}$ линеал


$\overset{\,\circ}{\cal F}$ такой, что при любых


$\lambda ,\mu _0,\mu _1,\dotsc,\mu _p\in


S_{a,\theta }^L(M)$, $f\in\overset{\,\circ}{\cal F}$


\begin{align*}


\|M(\lambda L-M)^{-1}L_{(\mu ,p)}^L(M)f\| _{{\cal F}}&\le


C(f)\Big/\bigg(|\lambda-a |\prod\limits_{k=0}^p|\mu _k-a|\bigg),\\


%


\|R_{(\mu ,p)}^L(M)(\lambda L-M)^{-1}\|_{\cal L(\cal F;\cal U)}&\le


K\Big/\bigg(|\lambda-a |\prod\limits_{k=0}^p|\mu_k-a|\bigg).


\end{align*}


\end{defn}








\begin{thm}


Пусть оператор $M$ сильно $(L,p)$-секториален. Тогда 


\begin{itemize}


\item[(i)] существует аналитическая в секторе


$\Sigma =\{t\in{\Bbb C}: |\arg t|<\theta-\pi/2\}$,


сильно непрерывная в нуле справа разрешающая полугруппа 


$\{ U^t:t\in\{0\}\cup\Sigma \}$ уравнения $(2);$ 


\item[(ii)] ${\cal U}={\cal U}^0\oplus{\cal U}^1$,


${\cal F}={\cal F}^0\oplus {\cal F}^1$, 


где ${\cal U}^0=\ker U^0$, ${\cal U}^1=\im U^0;$


\item[(iii)] $L_k\in\cal L(\cal U^k;\cal F^k)$,


$M_k\in \cal C l(\cal U^k;\cal F^k)$, $k=0,1;$


\item[(iv)] существуют операторы $M_0^{-1}\in\cal L(\cal F^0;\cal U^0)$, 


$L_1^{-1}\in \cal L(\cal F^1;\cal U^1)$, 


оператор $H=M_0^{-1}L_0\in\cal L(\cal U^0)$


нильпотентен степени не больше $p$.


\end{itemize}


\end{thm}





Через $M_k$ ($L_k$)


обозначено сужение оператора $M$ ($L$) на $\dom M_k={\cal


U}^k\cap\dom M$ (${\cal U}^k$), $k=0,1$.








\begin{note}


Ядра операторов полугруппы уравнения (2) в


условиях теоремы 1 совпадают с $M$-корневым линеалом оператора


$L$, состоящим из собственных и $M$-присоединенных высоты не


больше $p$ векторов оператора $L$. 


\end{note}








\section{Ограниченные на прямой решения}








Пусть оператор $M$ сильно $(L,p)$-секториален и $L$-спектр $\sigma^L(M)$ 


оператора $M$ не пересекается с мнимой осью. Обозначим


$$


\sigma_+=\{\mu\in\sigma^L(M):


\re\mu>0\},\quad\sigma_-=\{\mu\in\sigma^L(M):\re\mu<0\}.


$$


В силу замкнутости относительного спектра [1] существуют конечный


контур ${\Gamma}_+\subset\{\mu\in{\Bbb C}:\re\mu>0\}$,


ограничивающий $\sigma_+$, и неограниченная кривая ${\Gamma}_-


\subset\{\mu\in{\Bbb C}:\re\mu<0\}$, огибающая $\sigma_-$


таким образом, что при $\mu\in\Gamma_-$, $|\mu|\to\infty$ имеем


$\arg\mu\to\pm\theta$. В обоих случаях обход контура происходит


против часовой стрелки. Тогда существуют интегралы 


\begin{gather*}


R^t_+=\frac{1}{2\pi i}\int_{\Gamma_+ }(\mu L-M)^{-1}e^{\mu t}d\mu,


\quad t\in\Bbb R\,; \\


%


R^t_-=\frac{1}{2\pi i}\int_{\Gamma_- }(\mu L-M)^{-1}e^{\mu t}d\mu,


\quad t>0.


\end{gather*}








\begin{defn}


Оператор-функцию


$$


G^{t} =\begin{cases}


-R_+^{t}, & t<0; \\


\phantom{-}R_-^{t}, & t>0,


\end{cases}


$$


назовем {\it функцией Грина} уравнения (2).


\end{defn}





Используя относительно спектральную теорему ([11]; см. также [12]), получим 


следующие результаты. (Символом $Q$ обозначен проектор вдоль ${\cal F}^0$ на


${\cal F}^1$, существующий в силу теоремы~1.)








\begin{lemnonum}


Пусть оператор $M$ сильно


$(L,p)$-секториален и $L$-спектр оператора $M$ не пересекается с


мнимой осью. Тогда


\begin{itemize}


\item[(i)] $G^{t} : {\cal {F}} \to \dom M_{1};$


\item[(ii)] при $t\in{\Bbb R}\setminus\{0\}$ функция $G^{t}$ непрерывно 


дифференцируема и удовлетворяет


уравнению $L\frac{dG^{t}}{dt}=MG^{t}$, кроме того,


$$


\frac {dG^{t}}{dt} :


{\cal {F}} \to {\cal {U}}^{1}\quad


\text{и}


\quad\exists


C>0 \ \ \forall t\in{\Bbb R}\setminus\{0\}\quad \bigg\|\frac


{dG^{t}}{dt}\bigg\|_{\cal L(\cal F;\cal U)}\le \frac{C}{t}e^{at};


$$


\item[(iii)] $G^{0+}-G^{0-}=L_1^{-1}Q$.


\end{itemize}


\end{lemnonum}





Функцию $f:J\to{\cal F}$, где $J\subset{\Bbb R}\,$,


будем называть ограниченной,


если $\sup\limits_{t\in J}\| f(t)\|_{\cal F}<\infty$.


Пусть $k\in{\Bbb N}$, $Q\in{\cal L(\cal F)}$ --- проектор.


Обозначим через $C^{k,H}(J,{\cal


F};Q)$ класс функций $f$ таких,


что $f^0=(I-Q)f\in C^k(J,{\cal F})$, $f^1=Qf:J\to{\cal F}$


локально гельдерова. Символом $BC^k(J,{\cal F})$ обозначим множество 


функций $f\in C^k(J,{\cal F})$, для которых $f,f^{(1)},\dotsc,


f^{(k-1)}:J\to{\cal F}$ --- ограниченные функции. И, наконец, через


$BC^{k,H}(J,{\cal F};Q)$ обозначим множество таких функций $f$, что\linebreak


$(I-Q)f\in BC^k(J,{\cal F})$, 


$Qf:J\to{\cal F}$ --- локально гельдерова ограниченная функция.








\begin{thm}


Пусть оператор $M$ сильно


$(L,p)$-секториален и $L$-спектр оператора $M$ не пересекается с


мнимой осью. Тогда для любой функции $f \in BC^{p+1,H}({\Bbb


R},{\cal F};Q)$ уравнение $(2)$ имеет единственное


решение $u\in BC^{1}({\Bbb R},{\cal U}){:}$


\begin{equation}


u(t)= \int_{-\infty}^{+\infty} G^{t-s}f(s) ds -


\sum_{q=0}^p H^qM_0^{-1}f^{0(q)}(t).


\end{equation}


Если к тому же начальное значение $u_{0} \in {\cal U}$ имеет вид


$$


u_{0}= \int_{-\infty}^{+\infty} G^{-s}f(s) ds -


\sum_{q=0}^p H^qM_0^{-1}f^{0(q)}(0),


$$


то функция $(3)$ является единственным ограниченным на ${\Bbb R}$ решением


задачи $(1)$, $(2)$.


\end{thm}








\begin{thm}


Пусть оператор $M$ сильно


$(L,p)$-cекториален и для каждой функции $f \in BC^{p+1,H}({\Bbb


R},{\cal F};Q)$ существует единственное решение $u \in BC^1({\Bbb


R},{\cal U})$ уравнения $(2)$. Тогда\linebreak $L$-спектр оператора


$M$ не пересекается с мнимой осью. 


\end{thm}








\begin{note}


Если оператор $-M$ сильно $(L,p)$-секториален,


то утверждения теорем 2, 3 остаются в силе.


Действительно, функция $u(t)$ является решением уравнения (2)


тогда и только тогда, когда функция $v(t)=u(-t)$ разрешает уравнение $L\Dot


v(t)=-Mv(t)-f(-t)$.


\end{note}








\begin{thebibliography}{99}





\bibitem{1}


Свиридюк Г.А. {\it К общей теории полугрупп операторов} //


УМН. -- 1994. -- Т.\,49. -- Вып.\,4. -- С.\,47--74.


\bibitem{2}


Демиденко Г.В., Успенский С.В. {\it Уравнения и системы, не


разрешенные относительно старшей производной}. -- Новосибирск:


Научная книга, 1998. -- 438~с.


\bibitem{3}


Favini A., Yagi A. {\it Degenerate differential equations in


Banach spaces}. -- New~York--Basel--Hong Kong: Marcel Dekker,


Inc., 1999. -- 313~p.


\bibitem{4}


Егоров И.Е., Пятков С.Г., Попов С.В. {\it Неклассические


дифференциально-операторные уравнения}. -- Новосибирск: Наука,


2000. -- 336~c.


\bibitem{5}


Далецкий Ю.Л., Крейн М.Г. {\it Устойчивость решений


дифференциальных уравнений в банаховом пространстве}. -- М.: Наука,


1970. -- 534~с.


\bibitem{6}


Массера Х.Л., Шеффер Х.Х. {\it Линейные дифференциальные


уравнения и функциональные пространства}. -- М.: Мир, 1970. -- 456~с.


\bibitem{7}


Левитан Б.М., Жиков В.В. {\it Почти-периодические функции и


дифференциальные уравнения}. -- М.: Изд-во МГУ, 1978.


-- 205~c.


\bibitem{8}


Хенри Д. {\it Геометрическая теория полулинейных параболических


уравнений}. -- М.: Мир, 1985. -- 376~c.


\bibitem{9}


Келлер А.В. {\it Исследование ограниченных решений линейных


уравнений типа Соболева}: Дис.~$\dotsc$ канд. физ.-матем. наук. --


Челябинск: Челябинск. гос. ун-т, 1997. -- 115~c.


\bibitem{10}


Свиридюк Г.А., Федоров В.Е. {\it О единицах аналитических


полугрупп операторов с ядрами} // Сиб. матем. журн. -- 1998. -- Т.\,39.


-- \No\,3. -- С.\,604--616.


\bibitem{11}


Руткас А.Г. {\it Задача Коши для уравнения}


$Ax'(t)+Bx(t)=f(t)$


// Дифференц. уравнения. -- 1975. -- Т.\,11. -- \No\,11. -- С.\,1996--2010.


\bibitem{12}


Келлер А.В. {\it Относительно спектральная теорема} //


Вестн. Челяб. гос. ун-та. Сер. матем., мех. -- 1996. --


\No\,1. -- С.\,62--66.





\end{thebibliography}





\vskip 2cm





\postup{\it Челябинский государственный}{\qquad \it Поступили}


\postup{\it университет}{\it полный текст $31.07.2000$}


\postup{}{\it краткое сообщение $22.04.2004$}





\label{end}


\end{document}


